\documentclass[../../main.tex]{subfiles}

\begin{document}

\section{Stone-Weierstrass 定理}\label{section:ch12-Stone-Weierstrass}

以下定理是经典分析的瑰宝之一.

\begin{theorem}[Weierstrass 逼近定理]\label{theorem:ch12-Weierstrass}
令 $f$ 为闭有界区间 $[a,b]$ 上的连续实值函数. 则对每个 $\varepsilon>0$, 存在多项式 $p$ 满足
\begin{align*}
\text{对所有 }x\in[a,b],\quad |f(x)-p(x)|<\varepsilon.
\end{align*}
\end{theorem}

本节我们证明该定理的一个深远的推广. 对于紧 Hausdorff 空间 $X$, 考虑 $X$ 上具有最大值范数的连续实值函数空间 $C(X)$. Weierstrass 逼近定理告诉我们多项式在 $C[a,b]$ 中稠密.

现在 $C(X)$ 有了不是所有线性空间都具有的乘积结构, 即 $C(X)$ 中的两个函数 $f$ 与 $g$ 的乘积 $fg$ 仍然在 $C(X)$ 中. $C(X)$ 的线性子空间 $\mathcal{A}$ 称为一个代数, 若 $\mathcal{A}$ 中的任何两个函数的乘积也属于 $\mathcal{A}$. $X$ 上的实值函数族 $\mathcal{A}$ 称为分离 $X$ 中的点, 若对 $X$ 中的任何两个不同的点 $u$ 和 $v$, 存在 $\mathcal{A}$ 中的 $f$ 使得 $f(u)\neq f(v)$. 观察到由于 $X$ 是紧与 Hausdorff 的, 根据前一章的定理 18, 它是正规的, 因此我们可以从 Urysohn 引理推出整个代数 $C(X)$ 分离 $X$ 中的点.

\begin{theorem}[Stone-Weierstrass 逼近定理]\label{theorem:ch12-Stone-Weierstrass}
令 $X$ 为紧 Hausdorff 空间. 假定 $\mathcal{A}$ 是 $X$ 上连续实值函数的代数, 它分离 $X$ 中的点且包含常数函数. 则 $\mathcal{A}$ 在 $C(X)$ 中稠密.
\end{theorem}

观察到这是对 Weierstrass 逼近定理的推广, 因为闭有界区间 $[a,b]$ 是紧与 Hausdorff 的空间且多项式族是包含常数与分离点的代数.

在证明定理之前, 关于证明策略的一些说明是妥当的. 由于 $X$ 是紧与 Hausdorff 的, 它是正规的. 我们从 Urysohn 引理可以推断出对 $X$ 中的每对不交闭子集 $A$ 与 $B$ 以及 $\varepsilon\in(0,1/2)$, 存在函数 $f\in C(X)$ 使得
\begin{align*}
\text{在 }A\text{ 上 }f=\varepsilon/2,\quad\text{在 }B\text{ 上 }f=1-\varepsilon/2,\quad\text{且在 }X\text{ 上 }\varepsilon/2\leqslant f\leqslant 1-\varepsilon/2.
\end{align*}
因此, 若在 $X$ 上 $|h-f|<\varepsilon/2$, 则
\begin{align}
\text{在 }A\text{ 上 }h<\varepsilon,\quad\text{在 }B\text{ 上 }h>1-\varepsilon,\quad\text{且在 }X\text{ 上 }0\leqslant h\leqslant 1.\label{align:ch12-SW策略}
\end{align}
证明将分两步进行. 首先, 我们证明对 $X$ 的每对不交闭子集 $A$ 与 $B$ 以及 $\varepsilon\in(0,1/2)$, 存在属于代数 $\mathcal{A}$ 的函数 $h$ 使得\eqref{align:ch12-SW策略}成立. 我们接着证明 $C(X)$ 的任何函数 $f$ 可被这样的 $h$ 的线性组合一致逼近.

\begin{lemma}\label{lemma:ch12-SW点逼近}
令 $X$ 为紧 Hausdorff 空间, 而 $\mathcal{A}$ 是 $X$ 上的分离点且是包含常数函数的连续函数的代数. 则对 $X$ 的每个闭子集 $F$ 以及属于 $X\sim F$ 的点 $x_{0}$, 存在 $x_{0}$ 的邻域 $\mathcal{U}$ 与 $F$ 不交且具有性质: 对每个 $\varepsilon>0$, 存在函数 $h\in\mathcal{A}$ 使得
\begin{align}
\text{在 }\mathcal{U}\text{ 上 }h<\varepsilon,\quad\text{在 }F\text{ 上 }h>1-\varepsilon,\quad\text{且在 }X\text{ 上 }0\leqslant h\leqslant 1.\label{align:ch12-SW点邻域}
\end{align}
\end{lemma}

\begin{proof}
我们首先断言对每个点 $y\in F$, 存在 $\mathcal{A}$ 中的函数 $g_{y}$ 使得
\begin{align}
g_{y}(x_{0})=0,\quad g_{y}(y)>0\text{ 且在 }X\text{ 上 }0\leqslant g_{y}\leqslant 1.\label{align:ch12-SWgy}
\end{align}
事实上, 由于 $\mathcal{A}$ 分离点, 存在函数 $f\in\mathcal{A}$ 使得 $f(x_{0})\neq f(y)$. 函数
\begin{align*}
g_{y}=\left[\frac{f-f(x_{0})}{\|f-f(x_{0})\|_{\max}}\right]^{2}
\end{align*}
属于 $\mathcal{A}$ 且满足\eqref{align:ch12-SWgy}. 由 $g_{y}$ 的连续性, 存在 $y$ 的邻域 $\mathcal{N}_{y}$ 使得 $g_{y}$ 在其上仅取正值. 然而, $F$ 是紧空间 $X$ 的闭子集, 从而 $F$ 自身是紧的. 因此我们可以选取这些邻域的有限族 $\{\mathcal{N}_{y_{1}},\ldots,\mathcal{N}_{y_{n}}\}$ 覆盖 $F$. 定义函数 $g$ 为
\begin{align*}
g=\frac{1}{n}\sum_{i=1}^{n}g_{y_{i}}.
\end{align*}
则
\begin{align}
g(x_{0})=0,\quad\text{在 }F\text{ 上 }g>0,\quad\text{且在 }X\text{ 上 }0\leqslant g\leqslant 1.\label{align:ch12-SWg}
\end{align}
但紧集上的连续函数取到最小值, 因此我们可以选取 $c>0$ 使得在 $F$ 上 $g\geqslant c$, 通过可能地用正数乘以 $g$, 我们可以假定 $c<1$. 另一方面, $g$ 在 $x_{0}$ 是连续的, 因此存在 $x_{0}$ 的邻域 $\mathcal{U}$ 使得在 $\mathcal{U}$ 上 $g<c/2$. 因此 $g$ 属于代数 $\mathcal{A}$ 且
\begin{align}
\text{在 }\mathcal{U}\text{ 上 }g<c/2,\quad\text{在 }F\text{ 上 }g\geqslant c,\quad\text{而在 }X\text{ 上 }0\leqslant g\leqslant 1.\label{align:ch12-SWgc}
\end{align}
我们宣称\eqref{align:ch12-SW点邻域}对邻域 $\mathcal{U}$ 的这个选择成立. 令 $\varepsilon>0$. 根据 Weierstrass 逼近定理, 我们能够找到多项式 $p$ 使得
\begin{align}
\text{在}[0,c/2]\text{ 上 }p<\varepsilon,\quad\text{在}[c,1]\text{ 上 }p>1-\varepsilon,\quad\text{而在}[0,1]\text{ 上 }0\leqslant p\leqslant 1.\label{align:ch12-SW多项式}
\end{align}
由于 $p$ 是多项式而 $f$ 属于代数 $\mathcal{A}$, 复合 $h=p\circ g$ 也属于 $\mathcal{A}$. 从\eqref{align:ch12-SWgc}与\eqref{align:ch12-SW多项式}我们得出\eqref{align:ch12-SW点邻域}成立.
\end{proof}

\begin{lemma}\label{lemma:ch12-SW分离}
令 $X$ 为紧 Hausdorff 空间, $\mathcal{A}$ 是 $X$ 上的分离点且是包含常数函数的连续函数的代数. 则对 $X$ 的每对不交闭子集 $A$ 与 $B$ 以及 $\varepsilon>0$, 存在函数 $h$ 使得
\begin{align}
\text{在 }A\text{ 上 }h<\varepsilon,\quad\text{在 }B\text{ 上 }h>1-\varepsilon,\quad\text{且在 }X\text{ 上 }0\leqslant h\leqslant 1.\label{align:ch12-SW分离}
\end{align}
\end{lemma}

\begin{proof}
由前一个引理(其中 $F=B$), 对每个点 $x\in A$, 存在 $x$ 的邻域 $\mathcal{N}_{x}$ 与 $B$ 不交且具有性质\eqref{align:ch12-SW点邻域}. 然而, $A$ 是紧的, 由于它是紧空间 $X$ 的闭子集, 因此存在覆盖 $A$ 的有限邻域族 $\{\mathcal{N}_{x_{1}},\ldots,\mathcal{N}_{x_{n}}\}$. 选取 $\varepsilon_{0}$ 使得 $0<\varepsilon_{0}<\varepsilon$ 且 $(1-\varepsilon_{0}/n)^{n}>1-\varepsilon$. 对 $1\leqslant i\leqslant n$, 由于 $\mathcal{N}_{x_{i}}$ 具有性质\eqref{align:ch12-SW点邻域}, 其中 $B=F$, 我们选取 $h_{i}\in\mathcal{A}$ 使得
\begin{align*}
\text{在 }\mathcal{N}_{x_{i}}\text{ 上 }h_{i}<\varepsilon_{0}/n,\quad\text{在 }B\text{ 上 }h_{i}>1-\varepsilon_{0}/n,\quad\text{而在 }X\text{ 上 }0\leqslant h_{i}\leqslant 1.
\end{align*}
在 $X$ 上定义
\begin{align*}
h=h_{1}\cdot h_{2}\cdots h_{n}.
\end{align*}
则 $h$ 属于代数 $\mathcal{A}$. 由于对每个 $i$, 在 $X$ 上 $0\leqslant h_{i}\leqslant 1$, 我们有在 $X$ 上 $0\leqslant h\leqslant 1$. 对每个 $i$, 在 $B$ 上 $h_{i}>1-\varepsilon_{0}/n$, 因此在 $B$ 上 $h\geqslant(1-\varepsilon_{0}/n)^{n}>1-\varepsilon$. 最后, 对 $A$ 中的每个点 $x$ 存在指标 $i$ 使得 $x$ 属于 $\mathcal{N}_{x_{i}}$. 因此 $h_{i}(x)<\varepsilon_{0}/n<\varepsilon$. 由于对其他指标 $j$, $0\leqslant h_{j}(x)\leqslant 1$, 我们得出 $h(x)<\varepsilon$.
\end{proof}

\begin{proof}
令 $f$ 属于 $C(X)$. 设 $c=\|f\|_{\max}$. 若我们能用 $\mathcal{A}$ 中的函数任意逼近函数
\begin{align*}
\frac{f+c}{\|f+c\|_{\max}},
\end{align*}
对 $f$ 我们也能这样做. 因此我们可以假设在 $X$ 上 $0\leqslant f\leqslant 1$. 令 $n>1$ 为自然数. 考虑 $[0,1]$ 的分为长度为 $1/n$ 的 $n$ 个区间的一致划分 $\{0,1/n,2/n,\cdots,(n-1)/n,1\}$. 固定 $j$, $1\leqslant j\leqslant n$. 定义
\begin{align*}
A_{j}=\{x\text{ 属于 }X\mid f(x)\leqslant(j-1)/n\}\quad\text{与}\quad B_{j}=\{x\text{ 属于 }X\mid f(x)\geqslant j/n\}.
\end{align*}
由于 $f$ 是连续的, $A_{j}$ 和 $B_{j}$ 都是 $X$ 的闭子集, 当然它们是不交的. 根据前一个引理, 其中 $A=A_{j}$, $B=B_{j}$, 而 $\varepsilon=1/n$, 存在代数 $\mathcal{A}$ 中的函数 $g_{j}$ 使得
\begin{align}
\text{若 }f(x)\leqslant\frac{j-1}{n},g_{j}(x)<\frac{1}{n},\text{ 若 }f(x)\geqslant j/n,g_{j}(x)>1-1/n;\text{ 在 }X\text{ 上 }0\leqslant g_{j}\leqslant 1.\label{align:ch12-SWgj}
\end{align}
定义
\begin{align*}
g=\frac{1}{n}\sum_{j=1}^{n}g_{j}.
\end{align*}
则 $g$ 属于 $\mathcal{A}$. 我们宣称
\begin{align}
\|f-g\|_{\max}<3/n.\label{align:ch12-SW逼近3n}
\end{align}
一旦我们证明了这个断言, 证明就完成了, 这是由于给定 $\varepsilon>0$, 我们简单地选取 $n$ 使得 $3/n<\varepsilon$ 且因此 $\|f-g\|_{\max}<\varepsilon$. 为证明\eqref{align:ch12-SW逼近3n}, 我们首先证明
\begin{align}
\text{若 }1\leqslant k\leqslant n\text{ 且 }f(x)\leqslant k/n,\quad\text{则 }g(x)\leqslant k/n+1/n.\label{align:ch12-SW上估计}
\end{align}
的确, 对 $j=k+1,\ldots,n$, 由于 $f(x)\leqslant k/n$, $f(x)\leqslant(j-1)/n$, 因此 $g_{j}(x)\leqslant 1/n$. 因此
\begin{align*}
\frac{1}{n}\sum_{j=k+1}^{n}g_{j}\leqslant(n-k)/n^{2}\leqslant 1/n.
\end{align*}
由于对所有 $j$, $g_{j}(x)\leqslant 1$,
\begin{align*}
g(x)=\frac{1}{n}\sum_{j=1}^{n}g_{j}=\frac{1}{n}\sum_{j=1}^{k}g_{j}+\frac{1}{n}\sum_{j=k+1}^{n}g_{j}\leqslant\frac{1}{n}\sum_{j=1}^{k}g_{j}+1/n\leqslant k/n+1/n.
\end{align*}
因此\eqref{align:ch12-SW上估计}成立. 类似的方法可证明:
\begin{align}
\text{若 }1\leqslant k\leqslant n\text{ 且 }(k-1)/n\leqslant f(x),\quad\text{则 }(k-1)/n-1/n\leqslant g(x).\label{align:ch12-SW下估计}
\end{align}
对 $x\in X$ 选取 $k$, $1\leqslant k\leqslant n$, 使得 $(k-1)/n\leqslant f(x)\leqslant k/n$. 从\eqref{align:ch12-SW上估计}和\eqref{align:ch12-SW下估计}我们推出 $|f(x)-g(x)|<3/n$.
\end{proof}

我们以 Stone-Weierstrass 定理与 Urysohn 引理的优雅推论结束本章.

\begin{theorem}[Borsuk 定理]\label{theorem:ch12-Borsuk}
令 $X$ 为紧 Hausdorff 拓扑空间. 则 $C(X)$ 是可分的当且仅当 $X$ 上的拓扑可度量化.
\end{theorem}

\begin{proof}
首先假设 $X$ 可度量化. 令 $\rho$ 为诱导 $X$ 上的拓扑的度量. 则作为一个紧度量空间, $X$ 是可分的. 选取 $X$ 的可数稠密子集 $\{x_{n}\}$. 对每个自然数 $n$, 定义 $f_{n}(x)=\rho(x,x_{n})$. 由于 $\rho$ 诱导拓扑, $f_{n}$ 是连续的. 我们从 $\{x_{n}\}$ 的稠密性推出 $\{f_{n}\}$ 分离 $X$ 中的点. 在 $X$ 上定义 $f_{0}\equiv 1$. 现在令 $\mathcal{A}$ 为 $f_{k}(0\leqslant k<\infty)$ 中的有限项, 具有实系数的多项式族. 则 $\mathcal{A}$ 是包含常数函数的代数, 且它分离 $X$ 中的点, 因为它包含每个 $f_{k}$. 根据 Stone-Weierstrass 定理, $\mathcal{A}$ 在 $C(X)$ 中稠密. 但 $\mathcal{A}$ 中具有有理系数的多项式 $f$ 构成的函数族是可数集且在 $\mathcal{A}$ 中稠密. 因此 $C(X)$ 是可分的.

反过来, 假定 $C(X)$ 是可分的. 令 $\{g_{n}\}$ 为 $C(X)$ 的可数稠密子集. 对每个自然数 $n$, 定义 $\mathcal{O}_{n}=\{x\in X\mid g_{n}(x)>1/2\}$. 则 $\{\mathcal{O}_{n}\}_{1\leqslant n<\infty}$ 是开集的可数族. 我们宣称每个开集是 $\{\mathcal{O}_{n}\}_{1\leqslant n<\infty}$ 的子族的并, 因此 $X$ 是第二可数的. 但 $X$ 是正规的, 由于它是紧与 Hausdorff 的. Urysohn 度量化定理告诉我们 $X$ 是可度量化的. 为证明第二可数性, 令点 $x$ 属于开集 $\mathcal{O}$. 由于 $X$ 是正规的, 存在开集 $\mathcal{U}$ 使得 $x\in\mathcal{U}\subseteq\overline{\mathcal{U}}\subseteq\mathcal{O}$, 根据 Urysohn 引理, 存在 $C(X)$ 中的函数 $g$ 使得在 $\mathcal{U}$ 上 $g(x)=1$ 而在 $X\sim\mathcal{O}$ 上 $g=0$. 根据在 $C(X)$ 中的稠密性, 存在自然数 $n$ 使得在 $X$ 上 $|g-g_{n}|<1/2$. 因此 $x\in\mathcal{O}_{n}\subseteq\mathcal{O}$. 这完成了证明.
\end{proof}

\begin{exercise}\label{exercise:ch12-24}
假定 $X$ 是拓扑空间使得存在 $X$ 上的连续实值函数族. 证明 $X$ 必须是 Hausdorff 的.
\end{exercise}

\begin{exercise}\label{exercise:ch12-25}
令 $X$ 为紧 Hausdorff 空间而 $\mathcal{A}\subseteq C(X)$ 是包含常数函数的代数. 证明 $\mathcal{A}$ 在 $C(X)$ 中稠密当且仅当 $\mathcal{A}$ 分离 $X$ 中的点.
\end{exercise}

\begin{exercise}\label{exercise:ch12-26}
令 $\mathcal{A}$ 为紧空间 $X$ 上包含常数函数的实值连续函数的代数. 令 $f\in C(X)$ 具有性质: 对某个常数 $c$ 与实数 $\alpha$, 函数 $\alpha(f+c)$ 属于 $\overline{\mathcal{A}}$. 证明 $f$ 也属于 $\overline{\mathcal{A}}$.
\end{exercise}

\begin{exercise}\label{exercise:ch12-27}
对 $f,g\in C[a,b]$, 证明 $f=g$ 当且仅当对所有 $n$, $\int_{a}^{b}x^{n}f(x)\,dx=\int_{a}^{b}x^{n}g(x)\,dx$.
\end{exercise}

\begin{exercise}\label{exercise:ch12-28}
对 $f\in C[a,b]$ 与 $\varepsilon>0$, 证明存在实数 $c_{0},c_{1},\cdots,c_{n}$, 使得
\begin{align*}
\text{对所有 }x\in[a,b],\quad\left|f(x)-c_{0}-\sum_{k=1}^{n}c_{k}\cdot e^{kx}\right|<\varepsilon.
\end{align*}
\end{exercise}

\begin{exercise}\label{exercise:ch12-29}
对 $f\in C[0,\pi]$ 与 $\varepsilon>0$, 证明存在实数 $c_{0},c_{1},\cdots,c_{n}$, 使得
\begin{align*}
\text{对所有 }x\in[0,\pi],\quad\left|f(x)-c_{0}-\sum_{k=1}^{n}c_{k}\cdot\cos kx\right|<\varepsilon.
\end{align*}
\end{exercise}

\begin{exercise}\label{exercise:ch12-30}
令 $f$ 为 $\mathbf{R}$ 上以 $2\pi$ 为周期的连续实值函数. 对 $\varepsilon>0$, 证明存在实数 $c_{0},a_{1},\cdots,a_{n},b_{1},\cdots,b_{n}$, 使得
\begin{align*}
\text{对所有 }x\in\mathbf{R},\quad\left|f(x)-c_{0}-\sum_{k=1}^{n}[a_{k}\cos kx+b_{k}\sin kx]\right|<\varepsilon.
\end{align*}
(提示: 周期的连续函数可等同于平面的单位圆上的连续函数, 而单位圆是紧的且关于平面的拓扑是 Hausdorff 的.)
\end{exercise}

\begin{exercise}\label{exercise:ch12-31}
令 $X$ 与 $Y$ 为紧 Hausdorff 空间而 $f$ 属于 $C(X\times Y)$. 证明对每个 $\varepsilon>0$, 存在 $C(X)$ 中的函数 $f_{1},\cdots,f_{n}$ 与 $C(Y)$ 中的 $g_{1},\cdots,g_{n}$ 使得
\begin{align*}
\text{对所有 }(x,y)\in X\times Y,\quad\left|f(x,y)-\sum_{k=1}^{n}f_{k}(x)\cdot g_{k}(y)\right|<\varepsilon.
\end{align*}
\end{exercise}

\begin{exercise}\label{exercise:ch12-32}
在 Stone-Weierstrass 定理的证明中不用 Weierstrass 逼近定理, 证明存在自然数 $m$ 与 $n$ 使得多项式 $p(x)=1-(1-x^{n})^{m}$ 满足\eqref{align:ch12-SW多项式}. (提示: 由于 $p(0)=0$, $p(1)=1$ 且在 $(0,1)$ 上 $p'>0$, 仅须选取 $m$ 与 $n$ 使得 $p(c/2)<\varepsilon$ 且 $p(c)>1-\varepsilon$.)
\end{exercise}

\begin{exercise}\label{exercise:ch12-33}
令 $\mathcal{A}$ 为紧 Hausdorff 空间 $X$ 上的分离 $X$ 的点的连续实值函数族. 证明 $X$ 上的每个连续实值函数可被 $\mathcal{A}$ 中有限个函数的多项式任意近地一致逼近.
\end{exercise}

\begin{exercise}\label{exercise:ch12-34}
令 $\mathcal{A}$ 为紧 Hausdorff 空间 $X$ 上的连续实值函数的代数. 证明 $\mathcal{A}$ 的闭包 $\overline{\mathcal{A}}$ 也是代数.
\end{exercise}

\begin{exercise}\label{exercise:ch12-35}
令 $\mathcal{A}$ 为紧 Hausdorff 空间 $X$ 上的分离点的连续实值函数的代数. 证明或者 $\overline{\mathcal{A}}=C(X)$ 或者存在点 $x_{0}\in X$ 使得 $\overline{\mathcal{A}}=\{f\in C(X)\mid f(x_{0})=0\}$. (提示: 若 $1\in\overline{\mathcal{A}}$, 我们就成功了. 此外, 若对每个 $x\in X$ 存在 $f\in\mathcal{A}$ 满足 $f(x)\neq 0$, 则存在 $g\in\mathcal{A}$, 它在 $X$ 上是正的, 而这蕴涵 $1\in\overline{\mathcal{A}}$.)
\end{exercise}

\begin{exercise}\label{exercise:ch12-36}
令 $X$ 为紧 Hausdorff 空间, 而 $\mathcal{A}$ 是 $X$ 上的分离点且是包含常数函数的连续实值函数的代数.
\begin{enumerate}[(1)]
\item 给定任何两个数 $a$ 和 $b$ 与点 $u,v\in X$, 证明存在 $\mathcal{A}$ 中的函数 $f$, 使得 $f(u)=a$ 且 $f(v)=b$.
\item 给定任何两个数 $a$ 和 $b$ 与 $X$ 的不交闭子集 $A$ 和 $B$, 是否存在 $\mathcal{A}$ 中的函数 $f$, 使得在 $A$ 上 $f=a$ 而在 $B$ 上 $f=b$?
\end{enumerate}
\end{exercise}

\end{document}